\chapter{Od triangulacji do topologii i geometrii rozmaitości}
\label{ch:tg-most-geometria}

Ten rozdział scala trzy perspektywy:
triangulację jako komórkowy model przestrzeni,
graf dualny jako zapis sąsiedztwa komórek oraz
operatory łańcuchowe jako algebraiczne narzędzia
wykrywania cykli. Nie zastępujemy nimi pełnej
geometrii różniczkowej, lecz budujemy formalne
przejście do następnych części książki.

\section{Graf dualny nie koduje całej triangulacji}
\label{sec:tg-limit-dual}

Graf dualny triangulacji powierzchni zamkniętej jest
trójregularny (por. sekcja \ref{sec:tg-dual-triangulation});
w przypadku sklejeń boków dopuszczamy w nim także
krawędzie równoległe i pętle, liczone ze zwykłymi
krotnościami przy obliczaniu stopnia. Odwrotna konstrukcja nie jest jednak
jednoznaczna: do odtworzenia triangulacji trzeba wiedzieć,
które konkretne boki trójkątów są utożsamiane oraz
w jakich kierunkach. Jest to analogiczne do rozróżnienia
samego abstrakcyjnego grafu i jego układu rotacji.

\begin{przyklad}[Dwie ściany, dwie powierzchnie]
Bierzemy dwa abstrakcyjne trójkąty i sklejamy ich
boki parami. Oznaczmy wierzchołki obu trójkątów przez $a,b,c$
oraz $a',b',c'$. Sklejamy odpowiadające sobie boki
liniowo, tak aby $a\sim a'$, $b\sim b'$ i $c\sim c'$.
Jeżeli orientujemy oba dyski przeciwnie względem
tego przyporządkowania, sklejenie odwraca orientację
ich brzegów. Otrzymujemy dwa dyski sklejone
homeomorficznie wzdłuż całego brzegu, czyli sferę
(identyfikujemy je z dwiema półsferami).
Ma ona trzy klasy wierzchołków, trzy klasy krawędzi
i dwie ściany, zatem $\chi=3-3+2=2$.
Istnieje również sklejenie tych samych dwóch trójkątów
w kwadracie $[0,1]^2$ podzielonym przekątną
od $(0,0)$ do $(1,1)$, z identyfikacjami
$(0,t)\sim(1,t)$ oraz $(t,0)\sim(t,1)$.
Każda para boków (dwie pary z brzegu kwadratu
i para przekątnych) łączy dwa różne trójkąty.
Iloraz kwadratu jest homeomorficzny z $S^1\times S^1$
za pomocą $(x,y)\mapsto(e^{2\pi ix},e^{2\pi iy})$:
mapa indukuje ciągłą bijekcję ze zwartego ilorazu
do przestrzeni Hausdorffa. Wtedy jest jedna klasa wierzchołków,
trzy klasy krawędzi i dwie ściany, więc
$\chi=1-3+2=0$. Otrzymujemy torus.
W obu przypadkach abstrakcyjny multigraf dualny
ma dwa wierzchołki i trzy równoległe krawędzie.
Różnica topologiczna nie jest widoczna w grafie
dualnym bez dodatkowych informacji o sklejaniu.
Trzeba podkreślić, że otrzymane identyfikacjami
kompleksy są obiektami typu $\Delta$, nie zaś
abstrakcyjnymi kompleksami symplicjalnymi z dwoma
różnymi trójkątami.
\end{przyklad}

\begin{stwierdzenie}[Rekonstrukcja z danych incydencji]
Niech $\mathcal T$ będzie skończoną, niepustą rodziną
domkniętych trójkątów, których boki rozpatrujemy jako
osobne wystąpienia (również gdy pochodzą z tego samego
trójkąta). Załóżmy, że wszystkie te wystąpienia
podzielono na pary \emph{różnych} boków i dla każdej
pary podano homeomorfizm sklejający, odwzorowujący
końce jednego boku na końce drugiego.
Nie zakładamy, że iloraz jest spójny.
Wówczas iloraz jest jednoznacznie określoną
(z dokładnością do homeomorfizmu) zwartą
przestrzenią Hausdorffa. Ogniwo klasy wierzchołków
rozumiemy jako graf złożony z odcinków odpowiadających
narożom trójkątów, których końce utożsamiamy zgodnie
z parami boków przy tych narożach; zachowujemy
wszystkie wystąpienia, nawet w jednym trójkącie.
Iloraz jest powierzchnią bez brzegu: przy podanych
założeniach ogniwo każdej klasy wierzchołków jest
automatycznie pojedynczym okręgiem.
Składowe spójne ilorazu odpowiadają dokładnie
składowym spójnym grafu dualnego sklejonych
trójkątów. W szczególności, jeżeli graf dualny jest
spójny, otrzymujemy spójną powierzchnię zamkniętą.
\end{stwierdzenie}

\begin{proof}
Niech $X$ będzie rozłączną sumą trójkątów,
a $\sim$ relacją równoważności generowaną
przez zadane sklejenia. Relacja jest sumą
przekątnej $\{(x,x):x\in X\}\subset X\times X$,
grafów homeomorfizmów
sparowanych boków i ich odwrotności oraz
skończonego zbioru dodatkowych par równoważnych
wierzchołków. Istotnie, punkt we wnętrzu boku
należy tylko do jednego wystąpienia boku, więc
jego klasa składa się dokładnie z dwóch punktów;
punkty we wnętrzach trójkątów mają klasy
jednoelementowe. Dłuższe łańcuchy identyfikacji
mogą dotyczyć jedynie wierzchołków.
Każdy wymieniony graf jest zwarty, zatem
$\sim$ jest domknięta w $X\times X$.
Iloraz zwartej przestrzeni Hausdorffa przez
domkniętą relację równoważności jest zwarty
i Hausdorffa. Oznaczmy go przez $Y$,
a mapę ilorazową przez $q:X\to Y$.

Wyjaśnijmy również drugi aksjomat przeliczalności.
Niech $\mathcal B$ będzie przeliczalną bazą $X$,
a $\mathcal U$ rodziną jej skończonych sum.
Ponieważ $q$ jest domknięta (obraz zbioru zwartego
w przestrzeni Hausdorffa jest domknięty),
zbiory $Y\setminus q(X\setminus U)$ dla
$U\in\mathcal U$ są otwarte.
Dla $y\in O\subset Y$, gdzie $O$ jest otwarty,
zwarte włókno $q^{-1}(y)$ można pokryć skończenie
wieloma elementami $\mathcal B$ zawartymi w $q^{-1}(O)$.
Ich suma $U$ daje
$y\in Y\setminus q(X\setminus U)\subset O$.
Mamy więc przeliczalną bazę $Y$.

Możemy sprowadzić dane sklejenia do liniowych.
Dla pary boków z homeomorfizmem $g:e\to e'$
niech $\ell:e\to e'$ będzie liniowym odwzorowaniem
o tych samych obrazach końców. Na $e$ wybieramy
identyczność, a na $e'$ homeomorfizm
$\ell\circ g^{-1}$. Oba ustalają własne końce
i spełniają $(\ell\circ g^{-1})\circ g=\ell$.
Ponieważ każdy bok należy do jednej pary,
te wybory określają na brzegu każdego trójkąta
homeomorfizm $h$ ustalający wierzchołki.
Przedłużamy go do wnętrza trójkąta promieniowo
od środka ciężkości $o$: punkt
$(1-t)o+tx$, $x$ na brzegu, $0<t\le1$,
przechodzi na $(1-t)o+th(x)$, a $o$ pozostaje
nieruchomy. Jest to homeomorfizm, którego odwrotność
otrzymujemy z $h^{-1}$. Ciągłość w środku wynika
z ograniczoności trójkąta. Suma tych homeomorfizmów
przenosi dokładnie klasy pierwotnej relacji na klasy
relacji liniowej; ona i jej odwrotność indukują
homeomorfizmy ilorazów. Dotyczy to także pary
dwóch boków tego samego trójkąta.

W modelu liniowym punkty we wnętrzach trójkątów
mają otoczenia dyskowe. Przy punkcie we wnętrzu
sklejonego boku wybieramy w obu jego wystąpieniach
małe półdyski, używając zgodnego parametru wzdłuż
boku; sklejenie ich średnic daje dysk.
Przy klasie wierzchołków wybieramy w każdym jej
narożu sektor określony warunkiem, że współrzędna
barycentryczna danego wierzchołka przekracza
$1-\varepsilon$, gdzie $0<\varepsilon<1/2$.
Sektory są rozłączne przed sklejeniem, a ich suma
dla danej klasy wierzchołków jest otwarta i
nasycona względem relacji sklejania. Istotnie,
na każdym sparowanym boku liniowa identyfikacja
zachowuje parametr $r=1-t_v$, gdzie $t_v$ jest
współrzędną barycentryczną wierzchołka danego naroża.
Dla $r>0$ pozostałe dwie współrzędne mają sumę $r$;
po podzieleniu ich przez $r$ otrzymujemy parametr
odcinka ogniwa. Sektor domknięty $r\le\varepsilon$
jest stożkiem nad tym odcinkiem. Sklejając takie
domknięte sektory, otrzymujemy zwarty iloraz,
a współrzędne stożkowe dają ciągłą bijekcję na
stożek nad ogniwem $L$ (w punkcie $r=0$ ciągłość
wynika z topologii stożka i zwartości $L$).
Ponieważ $L$ jest skończonym grafem, stożek jest
Hausdorffa, więc bijekcja jest homeomorfizmem.
Ograniczenie do $r<\varepsilon$ daje otwarte
otoczenie stożkowe w $Y$.
Przy $r=0$ wszystkie naroża danej klasy są jednym
punktem. Każdy koniec odcinka ogniwa
jest sparowany z dokładnie jednym innym końcem;
wierzchołki grafu $L$ mają więc stopień dwa,
licząc pętle z krotnością. Zatem $L$ jest skończoną
sumą rozłącznych okręgów. Wykażemy, że jest spójne.
Każda identyfikacja dwóch wierzchołków przez
sklejenie boków utożsamia odpowiednie końce
odcinków ich naroży w $L$. Dwa naroża leżące
w tej samej klasie wierzchołków można połączyć
skończonym ciągiem takich identyfikacji, ponieważ
relacja na wierzchołkach jest generowana wyłącznie
przez sklejenia boków. Odcinki odpowiadające
kolejnym narożom w tym ciągu mają wspólny koniec
w $L$. Każde dwa odcinki $L$ należą więc do tej
samej składowej spójnej. Ponieważ $L$ jest
niepuste i spójne, jest jednym okręgiem.
Jego otwarty stożek jest otwartym dyskiem.

W bardziej ogólnym kompleksie, w którym dopuszczono
dodatkowe identyfikacje wierzchołków, spójność ogniwa
nie musi być automatyczna. Jeżeli ogniwo jest sumą
$k$ okręgów,
lokalna homologia w klasie wierzchołków w stopniu
dwa jest izomorficzna z $\widetilde H_1(L;\mathbb R)
\cong\mathbb R^k$: wynika to z wycięcia do otoczenia
stożkowego, jego kontraktowalności oraz ciągu
dokładnego pary z usuniętym wierzchołkiem,
jak w dowodzie twierdzenia \ref{thm:tg-link-surface}.
Dla punktu powierzchni bez brzegu ta grupa ma
wymiar jeden. Stąd koniecznie $k=1$.

Na koniec każda składowa grafu dualnego odpowiada
sumie obrazów trójkątów, które łączą się przez
wspólne boki, a zatem tworzą zbiór spójny.
Dwie różne składowe grafu nie mają wspólnych
punktów w ilorazie: każde utożsamienie, również
wierzchołków, jest generowane przez sklejenia boków.
Takie zbiory są zwarte i rozłączne, a jest ich
skończenie wiele, więc są także otwarte w $Y$.
Są dokładnie składowymi spójnymi $Y$.
W szczególności przy spójnym grafie dualnym
otrzymujemy spójną powierzchnię zwartą bez brzegu.
\end{proof}

\section{Laplasjan komórkowy i dyskretna teoria Hodge'a}
\label{sec:tg-cellular-hodge}

\begin{uwaga}[Dane dualne nie wyznaczają topologii]
Sam graf dualny triangulacji nie wyznacza na ogół typu
homeomorficznego rozmaitości. Aby odtworzyć kompleks
w wymiarze co najmniej dwa, potrzebujemy informacji,
które ściany są utożsamiane, jak działają odwzorowania
sklejające oraz jakie są ogniwa niżej wymiarowych sympleksów.
W ogólnych kompleksach, zwłaszcza wyższego wymiaru,
warunki lokalne, takie jak sferyczność ogniw, nie wynikają
z samej listy par przylegających komórek. Szczególny
przypadek pełnego parowania boków trójkątów
opisany powyżej nie wymaga dodatkowego testu ogniw.
\end{uwaga}

\begin{definicja}[Kołańcuchy i różniczka komórkowa]
Niech $K$ będzie skończonym kompleksem
symplicjalnym wymiaru $d$ i wybierzmy orientację
każdego sympleksu niezależnie od pozostałych.
Niech $C^p(K;\mathbb R):=
\operatorname{Hom}(C_p(K;\mathbb R),\mathbb R)$.
Korzystamy ze zwykłego, nieuzupełnionego kompleksu
łańcuchowego: poza stopniami $0,\ldots,d$
przestrzenie łańcuchów i kołańcuchów są zerowe.
Operatory brzegowe $\partial_p$ wyznaczają
\emph{różniczki kobrzegowe}
$d_p:C^p\to C^{p+1}$ przez $d_p\varphi=\varphi\circ\partial_{p+1}$.
Z równości $\partial_{p+1}\partial_{p+2}=0$
wynika $d_{p+1}d_p=0$, ponieważ dla $\varphi\in C^p$
mamy $(d_{p+1}d_p\varphi)(c)=
\varphi(\partial_{p+1}\partial_{p+2}c)=0$.
Macierz $d_p$ jest transpozycją macierzy
$\partial_{p+1}$ w bazach dualnych do wybranych
baz łańcuchów: jej współczynnik w wierszu $j$
i kolumnie $i$ jest wartością $i$-tego funkcjonału
bazowego na brzegu $j$-tego sympleksu.
Symbol $T$ będzie oznaczał transpozycję macierzy,
a nie sprzężenie względem dowolnego iloczynu
skalarnego. Przy ustalonych dodatnio określonych iloczynach skalarnych
w przestrzeniach $C^p$ przez $d_p^*$ oznaczamy
operator sprzężony, czyli
$\langle d_pu,v\rangle_{p+1}=\langle u,d_p^*v\rangle_p$.
Jeżeli $D_p$ jest macierzą $d_p$, a $G_p$ oraz
$G_{p+1}$ są macierzami Grama tych iloczynów,
to z równości $u^TD_p^TG_{p+1}v
=u^TG_p[d_p^*]v$ dla wszystkich $u,v$
wynika
\[
 [d_p^*]=G_p^{-1}D_p^TG_{p+1}.
\]
W szczególności zwykła transpozycja daje macierz
sprzężenia w bazach ortonormalnych. \emph{Laplasjanem Hodge'a}
na $p$-kołańcuchach nazywamy
\[
\Delta_p:=d_{p-1}d_{p-1}^*+d_p^*d_p,
\]
przyjmując brakujące składniki na końcach
kompleksu za zero.
\end{definicja}

\begin{stwierdzenie}[Dodatniość i harmoniczność]
Dla każdego $p$ laplasjan $\Delta_p$ jest
samosprzężony i nieujemny oraz
\[
\ker\Delta_p=\ker d_p\cap\ker d_{p-1}^*.
\]
\end{stwierdzenie}

\begin{proof}
Dla dowolnego $\omega\in C^p$ mamy
\[
\langle\Delta_p\omega,\omega\rangle
=\|d_p\omega\|^2+\|d_{p-1}^*\omega\|^2.
\]
Każdy składnik jest nieujemny, więc operator jest
nieujemny; jest samosprzężony jako suma operatorów
postaci $AA^*$ oraz $B^*B$. Jeżeli $\Delta_p\omega=0$,
oba kwadraty norm muszą znikać, zatem
$d_p\omega=0$ i $d_{p-1}^*\omega=0$.
Odwrotnie oba równania implikują
$\Delta_p\omega=0$.
\end{proof}

\begin{twierdzenie}[Skończeniewymiarowe twierdzenie Hodge'a]
\label{thm:tg-finite-hodge}
Niech $K$ będzie skończonym kompleksem symplicjalnym,
a w każdej przestrzeni $C^p(K;\mathbb R)$ ustalmy
dodatnio określony iloczyn skalarny. Wtedy zachodzi ortogonalny rozkład
\[
C^p=\operatorname{im}d_{p-1}
\ \oplus^{\perp}\ \ker\Delta_p
\ \oplus^{\perp}\ \operatorname{im}d_p^*,
\]
a naturalne odwzorowanie
$\ker\Delta_p\to H^p(K;\mathbb R)
:=\ker d_p/\operatorname{im}d_{p-1}$,
$\omega\mapsto[\omega]$, jest izomorfizmem.
\end{twierdzenie}

\begin{proof}
Z tożsamości $d_pd_{p-1}=0$ dostajemy
$\langle d_{p-1}u,d_p^*v\rangle
=\langle d_pd_{p-1}u,v\rangle=0$.
Zatem $U:=\operatorname{im}d_{p-1}$ oraz
$W:=\operatorname{im}d_p^*$ są ortogonalne.
Dla $\omega\in C^p$ warunek $\omega\perp U$
jest równoważny
$\langle d_{p-1}^*\omega,u\rangle=0$
dla każdego $u\in C^{p-1}$, czyli
$d_{p-1}^*\omega=0$. Analogicznie
$\omega\perp W$ wtedy i tylko wtedy, gdy
$d_p\omega=0$. Stąd wspólne dopełnienie to
$U^\perp\cap W^\perp=\ker\Delta_p$.
Ponieważ $U$ i $W$ są ortogonalne,
$U+W=U\oplus^\perp W$ jest podprzestrzenią
domkniętą (skończony wymiar).
Rozkład względem tej podprzestrzeni i jej
dopełnienia daje żądaną sumę ortogonalną.

Odwzorowanie do kohomologii jest dobrze określone
i liniowe: $\ker\Delta_p\subseteq\ker d_p$, a
$\operatorname{im}d_{p-1}\subseteq\ker d_p$.
Jeżeli $\omega\in\ker\Delta_p$ reprezentuje zerową
klasę kohomologii, to $\omega=d_{p-1}u\in U$,
a równocześnie $\omega\in U^\perp$, zatem $\omega=0$.
To daje injektywność. Dla surjektywności wybierzmy
zamknięty kołańcuch $\eta\in\ker d_p$. Rozłóżmy go
jako $\eta=d_{p-1}u+h+d_p^*v$.
Po zastosowaniu $d_p$ dostajemy $d_pd_p^*v=0$,
a zatem $\|d_p^*v\|^2
=\langle d_pd_p^*v,v\rangle=0$.
Stąd $\eta-d_{p-1}u=h\in\ker\Delta_p$,
co oznacza, że klasa $[\eta]$ ma reprezentanta
harmonicznego. Dowód jest zakończony.
\end{proof}

\begin{przyklad}[Trójkąt, jego wypełnienie i wymiar harmoniki]
Traktujmy graf cykliczny $C_3$ jako jednowymiarowy
kompleks symplicjalny, z orientacją krawędzi zgodną
z obiegiem cyklu. Przy standardowych iloczynach skalarnych
mamy $d_0=B^T$, $d_1=0$ i $\Delta_1=B^TB$,
zgodnie z rozdziałem \ref{ch:tg-laplacian}.
Dla kolejności krawędzi $a\to b$, $b\to c$,
$c\to a$ macierz incydencji i laplasjan wynoszą
\[
 B=\begin{pmatrix}-1&0&1\\1&-1&0\\0&1&-1\end{pmatrix},
 \qquad B^TB=\begin{pmatrix}2&-1&-1\\-1&2&-1\\-1&-1&2\end{pmatrix}.
\]
Ponieważ $B^TBx=0$ wtedy i tylko wtedy, gdy
$Bx=0$ (mnożymy z lewej przez $x^T$),
jądro jest rozpięte przez $(1,1,1)^T$.
Z twierdzenia \ref{thm:tg-finite-hodge} wynika
$\dim H^1(C_3;\mathbb R)=1$.

Po wypełnieniu cyklu jednym sympleksem dwuwymiarowym
$d_1:C^1\to C^2$ przypisuje kołańcuchowi sumę jego
trzech wartości na zorientowanym brzegu trójkąta.
Jest to niezerowe odwzorowanie na jednowymiarową
przestrzeń $C^2$, więc $\dim\ker d_1=2$.
Z twierdzenia \ref{thm:tg-rank} mamy także
$\dim\operatorname{im}d_0=2$.
Ponieważ $d_1d_0=0$, te podprzestrzenie są równe.
Stąd $H^1(\Delta^2;\mathbb R)=0$.
Przy tej orientacji trójkąta i standardowych
iloczynach skalarnych $d_1=(1,1,1)$, zatem
\[
 \Delta_1=B^TB+d_1^Td_1=3I_3.
\]
Rzeczywiście $\ker\Delta_1=0$ dla wypełnionego
trójkąta, zgodnie z twierdzeniem Hodge'a.
Rozróżnienie następuje na poziomie $d_1$,
nie na poziomie samego grafu wierzchołków i krawędzi.
\end{przyklad}

\section{Co będzie potrzebne w topologii trój- i czterowymiarowej}
\label{sec:tg-roadmap-topology}

\begin{uwaga}[Grafy diagramów, rozkładów i plumbing]
W kolejnym etapie przejdziemy od kombinatoryki
kompleksów do struktur bardziej geometrycznych.
Grafy Cayleya i grafy grup pozwolą kodować działania
i rozkłady grup fundamentalnych. Grafy Taita
przypisane diagramom węzłów pozwolą przenieść
niektóre rachunki z teorii węzłów na wielomiany grafowe.
Grafy rozkładów JSJ opisują sposób sklejania części
trójrozmaitości. Wreszcie ważone grafy plumbingowe
prowadzą do czterorozmaitości, których formy przecięcia
można odczytać z macierzy związanej z grafem.
Na razie istotne jest rozróżnienie trzech rodzajów danych:
(1) abstrakcyjnego grafu, (2) zanurzenia i danych sklejania,
(3) dodatkowych wag, orientacji i ramek geometrycznych.
Bez tego rozróżnienia późniejsze wzory na formy przecięcia
byłyby matematycznie nieprecyzyjne.
\end{uwaga}

\section*{Wskazówki bibliograficzne}
Standardowe wprowadzenie do kompleksów symplicjalnych i
homologii zawiera A. Hatcher, \emph{Algebraic Topology}.
Ruchy bistellarne i triangulacje rozmaitości omawia
U. Pachner, \emph{PL Homeomorphic Manifolds are Equivalent
by Elementary Shellings}.
Dyskretna teoria Morse'a jest przedstawiona w pracy
R. Formana, \emph{Morse Theory for Cell Complexes}.
Grafy Reeba i ich zastosowania do funkcji Morse'a
opisuje literatura topologii obliczeniowej, m.in.
H. Edelsbrunner i J. Harer, \emph{Computational Topology}.
Znane w skrypcie wyniki o charakterystyce Eulera,
homologii i rozkładach na uchwyty stanowią
przygotowanie do tych konstrukcji.


